\documentclass[11pt]{article}
\usepackage[margin=1in]{geometry}
\usepackage{amsmath,amssymb,microtype}
\usepackage[colorlinks=true,urlcolor=blue]{hyperref}
\title{Literature status: the compact-symmetry spectral-inequality problem}
\author{Automated mathematics run \texttt{fa\_banach\_001}, lane 2}
\date{11 August 2026}
\begin{document}
\maketitle

\noindent\textbf{Status:} literature already answered, affirmative for open
problem 2 under global matrix-symbol positivity.

\section*{Source question}
Duván Cardona, \emph{Spectral inequalities for elliptic
pseudo-differential operators on closed manifolds}, arXiv:2209.10690,
Section 3, physical page 18, second bullet.  The paper asks for which
manifolds the coordinate-invariance restriction $\rho\geq1-\delta$ can be
removed from Theorem 1.4 and conjectures that, on a manifold with symmetries
and under a suitable positivity condition on the symbol, its spectral
inequality remains valid for every $0\leq\delta<\rho\leq1$.

\section*{Supporting answer}
Duván Cardona, Julio Delgado, and Michael Ruzhansky,
\emph{Estimates for sums of eigenfunctions of elliptic pseudo-differential
operators on compact Lie groups}, arXiv:2209.12092.  On physical page 2 the
introduction states that validity in the complete range
$0\leq\delta<\rho\leq1$ was an open problem prior to that work.  Theorem 1.1,
starting on physical page 3, takes a compact Lie group $G$ and a positive
elliptic operator
\[
 A\in\Psi^m_{\rho,\delta}(G\times\widehat G),\qquad
 0\leq\delta<\rho\leq1,
\]
assumes its global matrix-valued symbol satisfies
$\sigma_A(x,\xi)\geq0$, and proves for every nonempty open
$\omega\subset G$ that
\[
 \|v\|_{L^2(G)}\leq C_1e^{C_2\lambda}\|v\|_{L^2(\omega)}
 \quad
 (v\in\operatorname{span}\{e_j:\lambda_j\leq\lambda\}).
\]
This is precisely the compact-symmetry realization requested by the source,
with the previously unspecified positivity condition made explicit.

\section*{Provenance and scope}
The supporting authors knew the problem: the supporting paper shares the
source author, cites Cardona's source preprint, and explicitly calls the same
complete-range statement previously open.  Therefore the provenance is
\texttt{literature\_already\_answered}, not an agent-inferred implication.

Only the second source bullet is resolved.  The first question, on suitable
boundary problems and in particular higher polyharmonic powers with boundary,
is outside this identification.  The global matrix-symbol condition is also
stronger than positivity of the principal symbol alone.

\section*{Evidence}
The packet contains exact local compilations of the stored arXiv sources as
\texttt{source\_paper.pdf} and
\texttt{supporting\_paper\_2209.12092.pdf}.  The source statement is on
physical page 18; the decisive supporting introduction and theorem begin on
physical pages 2 and 3.

\end{document}

